% Seitenränder -----------------------------------------------------------------
\setlength{\topskip}{\ht\strutbox} % behebt Warnung von geometry
\geometry{
%    showframe,
    a4paper,
    left=20mm,
    right=20mm,
    top=25mm,
    bottom=35mm,
    includehead=false, % head wird nicht von seitenrändern beeinflusst
    includefoot=false, % foot wird nicht von seitenrändern beeinflusst
}

\usepackage[
    automark, % Kapitelangaben in Kopfzeile automatisch erstellen
    headsepline, % Trennlinie unter Kopfzeile
    footsepline, % Trennlinie über Fußzeile
    ilines % Trennlinie linksbündig ausrichten
]{scrpage2}
% TODO old {scrpage2}
% TODO new {scrlayer-scrpage}

% Kopf- und Fußzeilen ----------------------------------------------------------
\pagestyle{scrheadings}
%\renewcommand{\chapterpagestyle}{scrheadings}
\clearscrheadfoot

% Kopfzeile
\renewcommand{\headfont}{\normalfont} % Schriftform der Kopfzeile
\ihead{\large{\textsc{\paperTitle}}\\}
\chead{}
\ohead{}
\setlength{\headheight}{15mm} % Höhe der Kopfzeile
%\setheadwidth[0pt]{textwithmarginpar} % Kopfzeile über den Text hinaus verbreitern (falls Logo den Text überdeckt) 

% Fußzeile
\ifoot{\\\version}
\cfoot{\\\today}
\ofoot{\\Seite \pagemark}
% abstand vom inhalt zur "grundlinie" (Boden vom Text) und nicht obere Kante vom Text
% 1em ist daher ca. das minimum
\setlength{\footskip}{10mm}

% Zeilenabstände
\singlespacing % Zeilenabstand 1,5 Zeilen
\frenchspacing % erzeugt ein wenig mehr Platz hinter einem Punkt 


% Schusterjungen und Hurenkinder vermeiden
\clubpenalty = 10000
\widowpenalty = 10000
\displaywidowpenalty = 10000

% Quellcode-Ausgabe formatieren
\lstset{numbers=left, numberstyle=\tiny, numbersep=5pt, breaklines=true}
\lstset{emph={square}, emphstyle=\color{red}, emph={[2]root,base}, emphstyle={[2]\color{blue}}}

\counterwithout{footnote}{section} % Fußnoten fortlaufend durchnummerieren
\setcounter{tocdepth}{2} % im Inhaltsverzeichnis werden die Kapitel bis zum Level der paragraphen übernommen
\setcounter{secnumdepth}{4} % Kapitel bis zum Level der paragraphen werden nummeriert
% \subsubsectionlevel = level 3 

% Aufzählungen anpassen
\renewcommand{\labelenumi}{\arabic{enumi}.}
\renewcommand{\labelenumii}{\arabic{enumi}.\arabic{enumii}.}
\renewcommand{\labelenumiii}{\arabic{enumi}.\arabic{enumii}.\arabic{enumiii}}

\renewcommand{\autodot}{} % Remove all end-of-counter dots

% Färbung:
\definecolor{main}{RGB}{0, 63, 107}

\definecolor{head}{RGB}{173,216,230} % {0.4, 0.6, 0.8}
\definecolor{sub}{RGB}{209,238,238}
\definecolor{odd}{RGB}{255,255,234}
\definecolor{hig}{RGB}{255,255,224}



% Farb und SChema Infos für Code
\definecolor{DEFhellgelb}{rgb}{1,1,0.9}
\definecolor{DEFKeys}{rgb}{0,0,1}
\definecolor{DEFIdentifier}{rgb}{0,0,0}
\definecolor{DEFComments}{rgb}{0,0.5,0}
\definecolor{DEFString}{rgb}{1,0,0}

\definecolor{JAVAred}{rgb}{0.6,0,0}
\definecolor{JAVAgreen}{rgb}{0.25,0.5,0.35}
\definecolor{JAVApurple}{rgb}{0.5,0,0.35}
\definecolor{JAVAdocblue}{rgb}{0.25,0.35,0.75}

\definecolor{forestgreen}{RGB}{34,139,34}
\definecolor{orangered}{RGB}{239,134,64}
\definecolor{darkblue}{rgb}{0.0,0.0,0.6}
\definecolor{gray}{rgb}{0.4,0.4,0.4}

\definecolor{HTMLdarkgray}{rgb}{0.4, 0.4, 0.4}
\definecolor{HTMLocherCode}{rgb}{1, 0.5, 0} % #FF7F00 -> rgb(239, 169, 0)
\definecolor{HTMLblueCode}{rgb}{0, 0, 0.93} % #0000EE -> rgb(0, 0, 238)
\definecolor{HTMLgreenCode}{rgb}{0, 0.6, 0} % #009900 -> rgb(0, 153, 0)

\lstset{
    float=hbp,
    basicstyle=\footnotesize,
    identifierstyle=\color{DEFIdentifier},
    keywordstyle=\color{DEFKeys},
    stringstyle=\color{DEFString},
    commentstyle=\color{DEFComments},
    backgroundcolor=\color{white},
    columns=flexible,
    tabsize=2,
    frame=single,
    extendedchars=true,
    showspaces=false,
    showstringspaces=false,
    numbers=left,
    numberstyle=\tiny,
    breaklines=true,
    breakautoindent=true,
    captionpos=b,
}

\lstdefinestyle{Java}{
    language=Java,
    tabsize=4,
    tabsize=4,
    breakatwhitespace=false,
    resetmargins=true
%
    basicstyle=\footnotesize\ttfamily,
    keywordstyle=\color{JAVApurple}\bfseries,
    stringstyle=\color{JAVAred},
    commentstyle=\color{JAVAgreen},
%
    sensitive=true,
    morecomment=[l]{//},
    morecomment=[s][\color{JAVAdocblue}]{/**}{*/},
    morestring=[b]",
    morestring=[b]',
    morekeywords={
        public,void,package,import,throws,
        abstract,and,array,as,break,case,catch,cfunction,class,clone,const,
        continue,declare,default,do,else,elseif,enddeclare,endfor,endforeach,
        endif,endswitch,endwhile,extends,final,for,foreach,function,global,
        goto,if,implements,interface,instanceof,namespace,new,old_function,or,
        private,protected,public,static,switch,throw,try,use,var,while,xor
        die,echo,empty,exit,eval,include,include_once,isset,list,require,
        require_once,return,print,unset},
}

\lstdefinestyle{XML} {
    basicstyle=\footnotesize,
    columns=flexible,
    tabsize=2,
    frame=single,
    extendedchars=true,
    showspaces=false,
    showstringspaces=false,
    numbers=left,
    numberstyle=\tiny,
    breaklines=true,
    breakautoindent=true,
    captionpos=b,
%
    language=XML,
    extendedchars=true,
    breaklines=true,
    breakatwhitespace=true,
    emph={},
    emphstyle=\color{red},
    basicstyle=\ttfamily,
    columns=fullflexible,
    commentstyle=\color{gray}\upshape,
    morestring=[b]",
    morecomment=[s]{<?}{?>},
    morecomment=[s][\color{forestgreen}]{<!--}{-->},
    keywordstyle=\color{orangered},
    stringstyle=\ttfamily\color{black}\normalfont,
    tagstyle=\color{darkblue}\textbf,
    morekeywords={attribute,xmlns,version,type,release},
}

\lstdefinelanguage{CSS} {
    morekeywords={
        color,background,margin,padding,font,weight,display,position,
        top,left,right,bottom,list,style,border,size,white,space,min,width},
    sensitive=false,
    morecomment=[l]{//},
    morecomment=[s]{/*}{*/},
    morestring=[b]',
    morestring=[b]"
}

\lstdefinelanguage{JavaScript} {
    showstringspaces=false,
    keywordstyle=\ttfamily\bfseries\color{darkblue},
    identifierstyle=\ttfamily,
    stringstyle=\ttfamily\color{forestgreen},
    commentstyle=\color{gray},
    rulecolor=\color{gray},
    xleftmargin=5pt,
    xrightmargin=5pt,
    aboveskip=\bigskipamount,
    belowskip=\bigskipamount,
%
    morekeywords={
        break,const,continue,delete,do,while,export,for,in,function,
        if,else,import,in,instanceOf,label,let,new,return,switch,this,
        throw,try,catch,typeof,var,void,with,yield
    },
    sensitive=false,
    morecomment=[l]{//},
    morecomment=[s]{/*}{*/},
    morestring=[b]",
    morestring=[d]'
}